%%%%%%%%%%%%%%%%%%%%%%%%%%%%%%%%%%%%%%%%%%%%%%%%%%%%%%%%%%%%%%%%%%%%%%%%%%%%%%%%%%%%%%%%%%%%%%%%%%%%%%
%
%   Filename    : abstract-english.tex 
%
%   Description : This file will contain your abstract in english language.
%                 
%%%%%%%%%%%%%%%%%%%%%%%%%%%%%%%%%%%%%%%%%%%%%%%%%%%%%%%%%%%%%%%%%%%%%%%%%%%%%%%%%%%%%%%%%%%%%%%%%%%%%%

\begin{abstract}
% Consists of 150 to 250 words in a \textcolor{blue}{\textbf{single paragraph}}, that provides the reader with a summary of your research/study and what you intend to do. It should be informative enough to serve as a substitute for reading the thesis document itself. The Abstract should contain the following sentences:
% \begin{itemize}
%    \item \textbf{Motivation}. One to two sentence(s) describing the research domain or problem area.
%    \item \textbf{Problem}. One to two sentence(s) describing the research challenge or opportunity to be addressed in your study
%    \item \textbf{Contribution}. One to two sentence(s) describing your intended contribution or how you plan to address the research question 
%    \item \textbf{Methodology}. One to two sentence(s) summarizing the approach or the manner in which you will carry out your study
%    \item \textbf{Significant Results or Findings}. One to two sentence(s) describing your findings and result (This can be skipped during proposal stage) 
% \end{itemize}

% Original
% Short-form video (SFV) platforms like TikTok, YouTube Shorts, and Instagram Reels have transformed digital content creation. However, content creators face uncertainty in predicting audience engagement before publishing. Significantly, early engagement often determines visibility, yet newly uploaded videos can lack interaction data, making pre-publication decisions challenging. Prior research has focused on platform-centered models using multimodal video features—visual, audio, and textual—while largely neglecting creator-related and contextual factors. This study develops a creator-oriented engagement prediction model that integrates multimodal content signals with creator characteristics, such as follower count and posting history, and contextual factors, including posting time and audience activity patterns. The model aims to estimate potential engagement before posting, enabling data-informed creative decisions. By bridging content, creator, and context, this research contributes to human-centered predictive analytics, supports strategic resource allocation in the creator economy, and may foster a more diverse and sustainable digital content ecosystem.

% Short-form video (SFV) platforms like TikTok, YouTube Shorts, and Instagram Reels have transformed digital content creation. However, content creators face uncertainty in predicting audience engagement before publishing. Significantly, early engagement often determines visibility, yet newly uploaded videos can lack interaction data, making pre-publication decisions challenging. Prior research has focused on platform-centered models using multimodal video features—visual, audio, and textual—while largely neglecting creator-related and contextual factors. As a result, no existing model directly enables individual creators to estimate their own video's engagement potential before uploading. This study develops a creator-oriented engagement prediction model that integrates multimodal content signals with creator characteristics, such as follower count and posting history, and contextual factors, including posting time and audience activity patterns. The model aims to estimate potential engagement before posting, enabling data-informed creative decisions. To achieve this, the study constructs a dataset through voluntary data donation from TikTok creators, combining raw video files with official creator analytics. This study leverages a frozen ensemble of Large Multimodal Models—VideoLLaMA2 and Qwen2.5-VL—to jointly predict engagement by fusing multimodal content embeddings with structured creator and contextual metadata. By bridging content, creator, and context, this research contributes to human-centered predictive analytics, supports strategic resource allocation in the creator economy, and may foster a more diverse and sustainable digital content ecosystem.

Short-form video (SFV) platforms like TikTok have transformed digital content creation. However, content creators face uncertainty in predicting audience engagement before publishing. Significantly, early engagement often determines visibility, yet newly uploaded videos can lack interaction data, making pre-publication decisions challenging. Prior research has focused on platform-centered models using multimodal video features—visual, audio, and textual—while largely neglecting creator-related and contextual factors. As a result, no openly documented model comprehensively enables individual creators to estimate their own video's engagement potential before uploading. This study will develop a creator-oriented engagement prediction model that integrates multimodal content signals with creator characteristics, such as follower count and account age, and contextual factors, including the time and day of posting. The model aims to estimate potential engagement—operationalized as two early viewer-retention metrics, the Engagement Continuation Rate (ECR) and Normalized Average Watch Percentage (NAWP)—before posting, enabling data-informed creative decisions. To achieve this, the study will construct MicroTok-PH, a dataset that will be built through voluntary data donation from Filipino micro-influencers, combining raw video files with official creator analytics. The proposed model, C3-LMM, will leverage a frozen ensemble of Large Multimodal Models—VideoLLaMA2 and Qwen2.5-VL—to predict engagement as a regression task, through two fusion strategies the study compares: prompt-based feature injection and an auxiliary-model ensemble. By bridging content, creator, and context, this research will contribute to human-centered predictive analytics, support strategic resource allocation in the creator economy, and may foster a more diverse and sustainable digital content ecosystem.

%
%  Do not put citations or quotes in the abract.
%

% Keywords can be found at \url{http://www.acm.org/about/class/class/2012?pageIndex=0}.  Click the 
% link ``HTML'' in the paragraph that starts with ''The \textbf{full CCS classification tree}...''.

\begin{flushleft}
\begin{tabular}{lp{4.25in}}
\hspace{-0.5em}\textbf{Keywords:}\hspace{0.25em} & short-form video, engagement prediction, pre-publication analytics, multimodal learning, creator-oriented modeling
\end{tabular}
\end{flushleft}
\end{abstract}